\section{Participants and execution}\label{participants-and-the-two-kinds-of-work}

An \textit{A} miner trains a versioned text model for origin, share, span, and
attribution outputs. The artifact contains the weights, tokenizer, calibration,
and execution contract needed to run every output it declares. The miner commits only its
digest on chain and uploads the artifact privately to the coordinator, which
publishes it once the cohort ends. A published artifact, the miner's own
included, earns nothing when entered again, so a miner earns again only with a
model that differs measurably from every published one or improves on its
source by half the promotion margin (Section~\ref{public-copies}).
Slop Ninja executes
submitted artifacts for the pilot's hidden evaluations and published ones for
customer requests. Neither step requires the miner to operate a serving
endpoint or Cloudflare account.
Miners pay for their training and development; the operator pays for centrally
served \textit{A} inference. Benchmark rewards depend on measured performance,
not a claim about training expenditure.

Slop Ninja is the subnet \emph{operator}. The \emph{subnet owner} is the account
recognized on chain; Slop Ninja may operate through ownership or a lease.
A lease must secure an owner-authorized execution path for required
hyperparameter changes, including enabling commit--reveal, and any owner-only
burn operations. Operating the service alone grants neither those powers nor
the owner's emission entitlement.

One lead validator, operated by Slop Ninja, serves as the \emph{coordinator}.
This validator is the serverless service: it administers the daily frozen
evaluation cohorts, runs their models and the reference models, and
publishes a signed package of per-miner scores and the public inputs needed
to calculate rewards. It also calculates, signs, submits, and recovers its own
native weight commitments, then certifies its finalized row. These tasks all
run within the serverless deployment.
Every validator verifies the package and its receipt index, independently
applies the frozen reward rules, and derives the canonical
native weight rows. The same package and registration bindings must produce
identical integer UID/weight entries. Each validator uses its own key, nonce,
and recovery journal to commit and reveal those rows in parallel with the
others. No validator waits for the coordinator's reveal or finalized
row before committing. Held-out documents, labels, individual predictions, customer text,
and private model weights remain outside the public package. Other validators
do not reproduce the coordinator's hidden model scoring.
Section~\ref{cohorts} defines the package bindings, native submission rules,
and per-validator recovery requirements.

This pilot trusts the coordinator to execute and score honestly. Validator
recomputation checks reward arithmetic over the published scores; it cannot
establish that the reported predictions are real. After its own row finalizes,
the coordinator publishes a certificate for audit against the same
score package and derived row. Its certificate authenticates those bindings
without creating a validator quorum or a submission prerequisite.
The coordinator also decides when measured quality justifies offering paid
inference, and which workloads the service may accept. Commercial activation
does not revise a cohort's scores.

Dependence on other validators is confined to their timely native weight
delivery. Scoring, cohort advancement, and customer service require no validator
quorum and do not depend on other validators' acknowledgments, availability,
software updates, or migration completion. Each validator submits and recovers
independently; delayed validator rows affect native reward allocation
without holding evaluation or service work open. The coordinator's own chain
delivery and certificate publication also do not gate cohorts or other validators.

Curators maintain authorized evaluation collections, production histories,
hidden labels, and source-group splits. Access to these materials follows the
collection's permissions. Customers authorize processing of submitted documents
and reference writing under their declared rights for a quoted job; buying inference gives no automatic
permission to add either to training or benchmarks. Author references and
derived profiles receive the same content protection as the query document.

Later \textit{C} contributors submit public grammar and lexical improvements.
Accepted contributions receive credit once under a separately activated
contribution contract (Section~\ref{grammar-contributions}). Later \textit{B}
miners register private transformation pipelines, with language-model adapters
served by Cloudflare in their own accounts and numeric rankers executed by
Slop Ninja. These miners maintain their delegated access, capacity, and billing,
and pay for their model calls, including required evaluations. Customer
payments follow the burn policy; no direct per-job reimbursement is promised.
Neither \textit{B} service obligations nor \textit{C} rewards apply to the
\textit{A}-only launch.

\subsection{Rust applications and serverless execution}

All Slop Ninja production application logic is written in Rust, including
native miner and non-coordinator validator clients, update and migration tools,
and the serverless coordinator, challenge, and inference logic. A \emph{challenge}
is a coordinator-internal evaluation task, such as generating cohort documents
or running an admitted artifact on them; no challenge is sent to a miner or
validator. Native participants use
\texttt{rtao} for chain integration~\cite{rtao}; Linux and macOS releases require
explicit platform qualification. Research ML backends remain separate from
this production contract. Native participants retain their own hotkeys.
The coordinator's validator hotkey is controlled within the serverless service,
with durable signing ownership and transaction recovery. Its authorization and
receipt keys remain distinct from that hotkey; other validators' keys never
enter the service.

Miner and validator clients install signed updates automatically and
independently. A participant that has missed several releases follows the
supported sequence of resumable schema migrations before resuming its work.
Updates preserve unfinished obligations and reconcile chain submissions before
retrying them. No participant waits for another to update or migrate. The
release, migration, and conditional rollback requirements appear in
Section~\ref{cohorts}; automatic recovery must satisfy its compatibility and
journal checks before new signing is allowed.

The coordinator is entirely serverless on Cloudflare Workers and Durable
Objects~\cite{cfworkers,cfdo}. Shared Rust code defines canonical messages,
artifact identities, scoring, and state transitions; supported modules compile
to Wasm with the required platform bindings and generated deployment glue.
The full native \texttt{rtao} runtime is not assumed to compile into a Worker.
Workers code follows a bounded-memory rule: it allocates only inside a fixed
arena sized from measured peaks, checks every declared length against that
budget before allocating, and streams large objects instead of buffering or
re-reading them.
The coordinator needs qualified Workers-compatible chain transport, encoding,
signing, and durable reconciliation, including its validator commitments and
reveal obligations. Both coordinator and native participants use the
\texttt{rtao-rpc} service, which provides Slop Ninja RPC
access at no service charge~\cite{rtaorpc}. Chain fees and execution costs remain
separately funded.

Live \texttt{rtao-rpc} webhooks, offered for both Finney and testnet, deliver
every finalized block. That per-block delivery is the coordinator's primary
heartbeat and clock, and each delivery is its own Worker invocation~\cite{cffetch}.
Durable Objects retain the finalized block cursor, and RPC reads backfill gaps
and resolve finality. A one-minute Cron Trigger~\cite{cfcron} is the only
fallback: when deliveries stall, including after the provider pauses a
subscription whose retries are exhausted, it processes blocks from public
OpenTensor RPC through the same per-block path. No permanently running
subscriber is required. A webhook is a notification, not evidence of finality or
permission to sign. Section~\ref{block-notifications} specifies ingress, the
cursor, fallback, and the operator-resumed pause, a stated limit; the webhook
receiver and the serverless chain path remain integration work.

\begin{table}[H]
\centering
\begin{tabularx}{\linewidth}{@{}lX@{}}
\toprule
Component & Responsibility \\
\midrule
API Workers & Enrollment, authenticated requests and block-webhook ingress,
status, and result retrieval. \\
Durable Objects & Authoritative cohort checkpoints, admission, reservations,
accounting, finalized-block cursors, and the coordinator's chain journal and
signing ownership. Each record has one transactional owner. \\
Lane Durable Objects & Execute fan-out lanes of bulk evaluation work. Each lane
claims tasks under the cohort object's durable leases, runs in its own
invocation, drives its own Container instance for model runs, and persists
outcomes before releasing its claim. \\
Queues & Carry lane wakeups holding identifiers and generations only; durable
lane records reconstruct any lost, duplicated, or expired message. \\
R2 storage & Versioned artifacts, evidence, manifests, and encrypted private
job objects, bound to their durable records by identity and digest. \\
Cron watchdog & Detect stalled deliveries or cursor lag, then process
finalized blocks from public RPC through the per-block path until webhooks
resume. \\
\bottomrule
\end{tabularx}
\caption{The serverless coordinator. Queue delivery, lane invocations, and
wakeups do not own protocol state.}
\end{table}

Task identities and dispatch intent are persisted before queue delivery, and a
lane's assignment is recorded in the cohort object before its wakeup is sent.
Messages carry identifiers and checkpoint generations, never customer content
or credentials. A lane checks the current authorization and claims the
recorded task before acting. Duplicate delivery cannot create a second
accounting outcome. External execution can still repeat after an uncertain
timeout; bounded retries retain every attempt and apply a declared selection
rule. Queue retention is finite, so durable records must reconstruct pending
dispatches after a long pause~\cite{cfqueues}. Cross-object work uses persisted
idempotent messages; a transaction in one Durable Object cannot atomically
update another~\cite{cfdo}.

All production inference, challenge generation and execution, and scoring run
serverlessly. Parsing and model execution take place outside the objects that
serialize admission and accounting. Bounded Rust/Wasm computation runs in
Workers; a qualified Cloudflare serverless Container under Slop Ninja's control
supplies bounded computation beyond their limits, including every \textit{A}
model run. Workers AI supplies supported language-model and adapter formats; it
cannot host the \textit{A} encoder
(Section~\ref{existing-evidence-and-implementation-status}). Bulk evaluation
work runs as fan-out jobs: the cohort's Durable Object divides challenge
generation, paired scoring, receipt bundling, and per-model score summaries
into lanes, at least one per scored artifact and more as the batch grows
(Section~\ref{fan-out}). Each scoring lane drives its own Container instance,
and generation lanes share one rate limiter. Lane count changes elapsed time and
per-lane start cost, never a document's outcome. Every workload needs measured CPU, memory, bundle,
and numerical-runtime limits before admission; these targets do not imply
support for arbitrary uploaded architectures. Miner-supplied executable plugins
are excluded. Coordinator state and scheduling remain in Workers and Durable
Objects, and production serving never depends on a miner's or validator's local
machine being online.

The native adapter uses durable \texttt{rtao} journal and weight/reveal recovery
APIs. It starts paused, reconciles finalized history, and requires explicit
resumption with application compatibility checks. This is an integration
component, not a deployed validator. Its tested scope and current limitations
are recorded in the ``Implementation boundary'' paragraph in
Section~\ref{native-implementation-boundary}; unsupported paths
require completion and qualification before deployment. The coordinator's
Durable Object journal must implement the same intent, attempt, and finalized
recovery requirements as the native validators' local journals; the native
file-backed implementation does not supply that storage adapter. The cohort
recovery contract in Section~\ref{cohorts} applies to both deployments, including
the evidence needed to resume after runtime upgrades or missed tempo boundaries.
